% Part: incompleteness
% Chapter: representability-in-q
% Section: undecidability

\documentclass[../../../include/open-logic-section]{subfiles}

\begin{document}

\olfileid{inc}{req}{und}
\olsection{ناپرېکړتيا}

تيوري~$\Th{T}$ هغه وخت \emph{ناپرېکړه‌وړ} بولو چې داسې محاسبوي
کړنلاره نۀ وي چې له متناهي شمېر ګامونو وروسته د~$\Th{T}$ د ژبې د
هرې جملې~$!A$ دپاره په باوري ډول دې پوښتنې ته سم ځواب ورکړي:
«آيا~$\Th{T}$، $!A$ ثابتوي؟» نو~$\Th{Q}$ به هله او يوازې هله د
پرېکړې وړ وه چې داسې محاسبوي کړنلاره وای چې د حساب د ژبې د ورکړل شوې
جملې~$!A$ دپاره پرېکړه وکړي چې~$\Th{Q} \Proves !A$ ده که نۀ۔ دا
پوښتنه داسې لا کره کولے شو: آيا اړيکه~$\Prov[\Th{Q}](y)$، چې پر~$y$
هله صدق کوي چې~$y$ په~$\Th{Q}$ کښې د ثابتېدونکې جملې ګوډل شمېره
وي، بازګشتي ده؟ ځواب «نه» دے۔

\begin{thm}
$\Th{Q}$ ناپرېکړه‌وړ ده؛ يعنې اړيکه
\[
\Prov[\Th{Q}](y) \defiff \fn{Sent}(y) \land
\lexists[x][\Prf[\Th{Q}](x, y)]
\]
بازګشتي نۀ ده۔
\end{thm}

\begin{proof}
فرض کړئ بازګشتي وای۔ بيا د درېدنې مسئله داسې حلولے شو: د ورکړو
$e$ او~$n$ دپاره پوهېږو چې~$\cfind{e}(n) \fdefined$ هله او يوازې
هله چې داسې~$s$ وي چې~$T(e, n, s)$، چې~$T$ د
\olref[cmp][rec][nft]{thm:kleene-nf} څخه د کلېني پريديکات دے۔ ځکه
$T$ بنسټيزه بازګشتي ده، په~$\Th{Q}$ کښې د کوم فارمول~$!B_T$ په
وسيله تمثيلېږي؛ يعنې~$\Th{Q} \Proves
!B_T(\num{e}, \num{n}, \num{s})$ هله او يوازې هله چې~$T(e,n,s)$۔
که~$\Th{Q} \Proves !B_T(\num{e}, \num{n}, \num{s})$، نو
$ \Th{Q} \Proves \lexists[y][!B_T(\num{e}, \num{n}, y)]$ هم۔ که
داسې~$s$ نۀ وي، نو د هر~$s$ دپاره
$\Th{Q} \Proves \lnot !B_T(\num{e}, \num{n}, \num{s})$۔ خو~$\Th{Q}$،
$\omega$-سازګاره ده؛ يعنې که د هر~$n \in \Nat$ دپاره
$\Th{Q} \Proves \lnot !A(\num{n})$، نو
$\Th{Q} \Proves/ \lexists[y][!A(y)]$۔ دا ځکه پوهېږو چې د~$\Th{Q}$
بديهي اصول په معياري مدل~$\Struct{N}$ کښې رښتوني دي۔ نو
$\Th{Q} \Proves/ \lexists[y][!B_T(\num{e}, \num{n}, y)]$۔ په بله
وينا، $\Th{Q} \Proves \lexists[y][!B_T(\num{e}, \num{n}, y)]$
هله او يوازې هله چې داسې~$s$ وي چې~$T(e,n,s)$؛ يعنې هله او يوازې
هله چې~$\cfind{e}(n) \fdefined$۔ له~$e$ او~$n$ څخه
$\Gn{\lexists[y][!B_T(\num{e}, \num{n}, y)]}$ محاسبه کولے شو؛
$g(e,n)$ هغه بنسټيزه بازګشتي تابعه ونيسئ چې دا کار کوي۔ نو
\[
h(e, n) =
\begin{cases}
1 & \text{که~$\Prov[\Th{Q}](g(e, n))$ وي}\\
0 & \text{که نۀ وي}.
\end{cases}
\]
دا به ښودلې وای چې که~$\Prov[\Th{Q}]$ بازګشتي وي، $h$ هم بازګشتي
ده۔ خو د~\olref[cmp][rec][hlt]{thm:halting-problem} له مخې~$h$
بازګشتي نۀ ده، نو~$\Prov[\Th{Q}]$ هم بازګشتي نۀ شي کېدے۔
\end{proof}

\begin{cor}
د لومړۍ درجې منطق ناپرېکړه‌وړ دے۔
\end{cor}

\begin{proof}
که د لومړۍ درجې منطق د پرېکړې وړ وای، په~$\Th{Q}$ کښې ثبوت‌وړتيا
هم د پرېکړې وړ وه، ځکه~$\Th{Q} \Proves !A$ هله او يوازې هله چې
$\Proves !T \lif !A$، چېرته چې~$!T$ د~$\Th{Q}$ د بديهي اصولو
عطف دے۔
\end{proof}

\end{document}
